% SSSUP-FONTE-STORICA: fonti-locali/originali/2014-2020.zip
% Anno 2019, esercizio 3. Trasposizione dal sorgente LaTeX storico.
% Verificare editorialmente la trasposizione prima di considerarla definitiva.
\renewcommand{\Anno}{2019}
\renewcommand{\NumeroProblema}{3}

\renewcommand{\Testo}{
Mostrare, preliminarmente, che tra tutti i triangoli di area $A$ il triangolo di perimetro $P$ minimo è quello equilatero. Dedurre la più piccola costante $C$ per cui vale la disuguaglianza
$$\sqrt{A}\le CP\qquad\hbox{per ogni triangolo}.$$
Incollare 2 triangoli vuol dire considerare la figura piana ottenuta unendo i 2 triangoli in modo che uno dei lati del primo triangolo coincida con uno del secondo (dunque i 2 triangoli devono avere almeno un lato uguale) e che l'interno delle figure non si intersechi.
\begin{enumerate}
\item[(a)] Mostrare che incollando 2 copie uguali o riflesse di uno stesso triangolo si ottiene un quadrilatero (anche non convesso) oppure un triangolo. In quali casi si ottiene un triangolo ed in quali un parallelogramma?
\item[(b)] Determinare il triangolo $T$ di area $A$ fissata tale che incollando 2 copie uguali o riflesse di $T$ si ottiene la figura di perimetro $P$ minimo. Determinare quale figura si ottiene, il suo perimetro, e la più piccola costante $C$ per cui
$$\sqrt{A}\le CP\qquad\hbox{per ogni figura ottenuta incollando due triangoli}.$$
\item[(c)] Svolgere il punto (b) sostituendo la parola triangolo con la parola quadrilatero, dove l'incollamento di due quadrilateri è definito in modo analogo.
\end{enumerate}
}

\renewcommand{\Risposta}{
% SSSUP-BOZZA-AI: Risposta; autore=ChatGPT; data=2026-11-12; stato=DA-VERIFICARE
% SSSUP-PROVENIENZA: ChatGPT (OpenAI), sessione ProblemiSSSUP del 2026-11-12; derivato dalla traccia canonica; revisione umana pendente.
Nel preliminare $C=1/\sqrt{12\sqrt3}$.

(a) L'incollamento che scambia gli estremi del lato comune produce sempre
un parallelogramma. Quello speculare produce un triangolo esattamente
quando il triangolo di partenza \`e rettangolo e si incolla lungo un cateto;
produce un parallelogramma esattamente quando il triangolo \`e isoscele
con base sul lato comune.

(b) Il pezzo ottimale \`e un triangolo rettangolo isoscele, incollato
lungo l'ipotenusa: si ottiene un quadrato e $P_{\min}=4\sqrt{2A}$.

(c) Il pezzo ottimale \`e met\`a di un esagono regolare tagliato lungo
una diagonale fra vertici opposti: un trapezio isoscele con basi $2a,a$
e lati obliqui $a$. Si ottiene l'esagono regolare e
$P_{\min}=4\sqrt[4]{3}\sqrt A$.

Se $A$ \`e l'area del singolo pezzo, le costanti ottimali in (b) e (c)
sono rispettivamente $1/(4\sqrt2)$ e $1/(4\sqrt[4]{3})$.
Se invece si usa l'area totale $S=2A$, le costanti nella disuguaglianza
$\sqrt S\le CP$ sono $1/4$ e $1/\sqrt{8\sqrt3}$.
}
\renewcommand{\Aiutino}{
% SSSUP-BOZZA-AI: Aiutino; autore=ChatGPT; data=2026-11-12; stato=DA-VERIFICARE
% SSSUP-PROVENIENZA: ChatGPT (OpenAI), sessione ProblemiSSSUP del 2026-11-12; derivato dalla traccia canonica; revisione umana pendente.
Per il preliminare, usare la formula di Erone.
Nell'incollamento il lato comune scompare dal bordo: il pezzo ottimale
non deve necessariamente essere quello di perimetro minimo preso da solo.
Per il punto (c), contare quanti vertici distinti pu\`o avere la figura
unita e confrontarla con il poligono regolare corrispondente.
Distinguere sempre l'area del singolo pezzo da quella della figura ottenuta.
}
\renewcommand{\SoluzioneUfficiale}{}
\renewcommand{\SoluzioneRedazionale}{
% SSSUP-BOZZA-AI: SoluzioneRedazionale; autore=ChatGPT; data=2026-11-12; stato=DA-VERIFICARE
% SSSUP-PROVENIENZA: ChatGPT (OpenAI), sessione ProblemiSSSUP del 2026-11-12; derivato dalla traccia canonica; revisione umana pendente.
\textbf{Notazione per le aree.}
Manteniamo $A$ per l'area di un singolo pezzo e indichiamo con $S=2A$
l'area della figura ottenuta dall'incollamento.
Il Testo usa $A$ prima per il pezzo e poi in una disuguaglianza riferita
alla figura: per evitare equivoci, riporteremo le costanti in entrambe
le convenzioni. I poligoni sono intesi semplici e non degeneri;
il perimetro della figura unita non comprende il lato divenuto interno.

\textbf{Preliminare: un solo triangolo.}
Siano $a,b,c$ i lati e $s=(a+b+c)/2$ il semiperimetro.
La formula di Erone e la disuguaglianza fra media aritmetica
e media geometrica danno
\[
 A^2=s(s-a)(s-b)(s-c)
 \le s\left(\frac{s}{3}\right)^3
 =\frac{s^4}{27}.
\]
L'uguaglianza si verifica esattamente quando $s-a=s-b=s-c$,
cio\`e quando il triangolo \`e equilatero.
Poich\'e $P=2s$, segue
\[
 A\le\frac{P^2}{12\sqrt3},\qquad
 P_{\min}=\sqrt{12\sqrt3\,A},\qquad
 C=\frac1{\sqrt{12\sqrt3}}.
\]
L'uguaglianza nel triangolo equilatero dimostra anche
che la costante non pu\`o essere ridotta.

\textbf{(a) Le forme possibili.}
Chiamiamo $U,V$ gli estremi del lato comune e $W$ il terzo vertice
del primo triangolo. La seconda copia deve stare nell'altro semipiano
rispetto alla retta $UV$. A meno di coincidenze dovute a simmetrie,
le possibilit\`a sono due.

Se gli estremi del lato si corrispondono senza scambiarsi, il nuovo
vertice $W'$ \`e il riflesso di $W$ rispetto alla retta $UV$.
La figura ha contorno $U,W,V,W'$. Indicando con $\alpha,\beta$
gli angoli del triangolo in $U,V$, i corrispondenti angoli dell'unione
sono $2\alpha,2\beta$. Il contorno \`e dunque un quadrilatero,
eventualmente concavo, salvo quando uno di questi angoli \`e piatto.
Questo accade esattamente se $\alpha=\pi/2$ oppure $\beta=\pi/2$:
il triangolo di partenza \`e rettangolo e il lato comune \`e un cateto.
In tal caso l'unione \`e un triangolo.

Nella stessa disposizione speculare, le diagonali del quadrilatero
sono $UV$ e $WW'$. Esse si tagliano a met\`a esattamente quando
$W$ appartiene all'asse di $UV$, cio\`e quando $UW=VW$.
Pertanto si ottiene un parallelogramma esattamente nel caso isoscele
con base $UV$; si tratta allora di un rombo.

Se invece si scambiano gli estremi, la seconda copia si ottiene
con una rotazione di $180$ gradi intorno al punto medio di $UV$.
Il nuovo vertice \`e $W''=U+V-W$, in notazione vettoriale.
Le diagonali $UV$ e $WW''$ hanno lo stesso punto medio,
quindi l'unione \`e sempre un parallelogramma non degenere.
Questo esaurisce le possibilit\`a anche quando alcuni lati del
triangolo sono uguali.

\textbf{(b) Due triangoli.}
Indichiamo con $a,b$ i due lati del singolo triangolo che non vengono
incollati e con $\gamma$ l'angolo compreso. Il perimetro esterno \`e
$P=2(a+b)$, mentre
\[
 A=\frac12 ab\sin\gamma
 \le\frac12 ab
 \le\frac{(a+b)^2}{8}
 =\frac{P^2}{32}.
\]
Ne segue $P\ge4\sqrt{2A}$.
Per avere uguaglianza occorrono e bastano $\gamma=\pi/2$ e $a=b$.
Il triangolo \`e dunque rettangolo isoscele; il lato incollato
\`e l'ipotenusa e l'unione \`e un quadrato di lato $\sqrt{2A}$.

Con l'area del singolo triangolo la costante ottimale \`e
$1/(4\sqrt2)$. Con l'area totale $S=2A$ la stessa conclusione
si scrive $\sqrt S\le P/4$, quindi la costante \`e $1/4$.

\textbf{(c) Due quadrilateri.}
Usiamo la seguente disuguaglianza poligonale, giustificata in fondo
alla soluzione: un poligono convesso con al pi\`u $m$ vertici,
area $B$ e perimetro $Q$, soddisfa
\[
 B\le\frac{Q^2}{4m\tan(\pi/m)},
\]
con uguaglianza soltanto per il poligono regolare con $m$ lati.

Due quadrilateri incollati lungo un lato hanno al pi\`u sei vertici
distinti complessivi. Il loro inviluppo convesso $H$ ha quindi
al pi\`u sei vertici. Passando all'inviluppo convesso, l'area
non diminuisce e il perimetro non aumenta: ogni rientranza viene
sostituita da una corda, pi\`u corta della corrispondente spezzata.
Indicando con $P_H$ il perimetro di $H$, otteniamo
\[
 2A=S\le\operatorname{area}(H)
 \le\frac{P_H^2}{24\tan(\pi/6)}
 =\frac{P_H^2}{8\sqrt3}
 \le\frac{P^2}{8\sqrt3}.
\]
Pertanto
\[
 P\ge4\sqrt[4]{3}\sqrt A.
\]
L'argomento copre anche quadrilateri concavi e ulteriori contatti
fra i bordi: per il limite inferiore conta soltanto l'inviluppo convesso.

Il limite si raggiunge prendendo un esagono regolare di lato $a$
e dividendolo con una diagonale fra vertici opposti.
Le due met\`a sono trapezi isosceli congruenti, con basi $2a,a$,
lati obliqui $a$ e altezza $\sqrt3\,a/2$.
L'area di ciascun trapezio \`e
\[
 A=\frac{(2a+a)}2\frac{\sqrt3}{2}a
 =\frac{3\sqrt3}{4}a^2,
\]
mentre il perimetro dell'esagono \`e $6a$.
Scegliendo $a=\sqrt{4A/(3\sqrt3)}$ si ottiene proprio
$P=4\sqrt[4]{3}\sqrt A$.

Inoltre, l'uguaglianza nelle disuguaglianze precedenti impone
che l'unione sia l'esagono regolare. Per dividerlo in due quadrilateri
con un lato in comune bisogna unire due vertici opposti;
il pezzo ottimale \`e dunque, a meno di similitudini e isometrie,
il trapezio appena descritto.

Le costanti ottimali sono perci\`o
\[
 \frac{\sqrt A}{P}\le\frac1{4\sqrt[4]{3}},
 \qquad
 \frac{\sqrt S}{P}\le\frac1{\sqrt{8\sqrt3}}.
\]

\textbf{Giustificazione della disuguaglianza poligonale.}
Fissiamo il perimetro $Q>0$ e massimizziamo l'area fra i poligoni
convessi con al pi\`u $m$ vertici. Un massimo esiste per compattezza,
dopo aver fissato un vertice per eliminare le traslazioni.
Un massimo ha esattamente $m$ vertici effettivi: se ne avesse meno,
si potrebbe sostituire un lato con due segmenti che formano un
piccolo triangolo esterno. Al variare dell'altezza $h$ di questo
triangolo, l'area aumenta di un termine proporzionale ad $h$,
mentre il perimetro aumenta solo di un termine proporzionale ad $h^2$.
Una riscalatura per ripristinare $Q$ lascerebbe quindi un aumento
di area per $h$ sufficientemente piccolo. Il poligono massimizzante
\`e dunque strettamente convesso.

Fissiamo ora i vicini di uno dei suoi vertici e facciamo variare
quel vertice mantenendo costante la somma dei due lati adiacenti.
Il luogo delle posizioni possibili \`e un'ellisse con fuochi nei vicini.
La parte variabile dell'area \`e l'area del triangolo formato dai tre
vertici. Il massimo locale si trova nel punto dell'ellisse pi\`u lontano
dalla retta dei fuochi, cio\`e sull'asse del segmento che li unisce.
I due lati adiacenti sono perci\`o uguali. Ripetendo il ragionamento,
tutti i lati del poligono sono uguali.

Per $m\ge4$, consideriamo quattro vertici consecutivi e variamo
i due intermedi, mantenendo fissi gli estremi e le lunghezze
dei tre lati consecutivi. Anche il quarto lato del quadrilatero,
che \`e una diagonale del poligono, resta fisso.
Se $t$ \`e una diagonale di questo quadrilatero e $\theta,\varphi$
sono i due angoli opposti a $t$, la derivata della sua area rispetto
a $t$ \`e
\[
 \frac{t}{2}\bigl(\cot\theta+\cot\varphi\bigr).
\]
Per un singolo triangolo la formula segue derivando il teorema
del coseno e l'espressione dell'area; le due derivate si sommano.
Le variazioni sono ammissibili localmente grazie alla stretta convessit\`a.
Al massimo la derivata \`e nulla, dunque
$\sin(\theta+\varphi)=0$ e quindi $\theta+\varphi=\pi$.
I quattro vertici sono perci\`o conciclici.
Applicando il ragionamento a gruppi consecutivi che hanno tre vertici
in comune, tutti i vertici del poligono appartengono alla stessa
circonferenza. Per $m=3$ questa propriet\`a \`e automatica.

Un poligono convesso equilatero e ciclico \`e regolare:
le corde uguali individuano archi consecutivi uguali. Non pu\`o
comparire un arco maggiore al posto di quello minore, perch\'e,
assieme agli altri $m-1$ archi, la somma supererebbe $2\pi$.
Se il lato \`e $Q/m$, l'apotema \`e
$Q/(2m\tan(\pi/m))$. L'area massima \`e pertanto
\[
 \frac12 Q\frac{Q}{2m\tan(\pi/m)}
 =\frac{Q^2}{4m\tan(\pi/m)},
\]
come richiesto, con uguaglianza precisamente nel caso regolare.
}
\renewcommand{\SoluzioniAlternative}{}
\renewcommand{\Osservazioni}{}
\renewcommand{\Commenti}{}
